\begin{textsample}{POS Dim 1 – pb – Score 63.00 – pb/pb\_em\_41.txt}  \label{ex:f1_pos_121}
2.1.6.4.2 \textbf{Fundamentos} teóricos/específicos do Teatro A \textbf{percepção} sensível do indivíduo, desde a instauração dos ditames da vida \textbf{moderna}, está comprimida por \textbf{uma} vivência urbana marcada pelo bombardeio de informações, propagandas e mídias sociais, \textbf{que} ditam tendências, comportamentos e desejos de consumo.
Muitas vezes o \textbf{que} \textbf{vemos} é o empobrecimento da experiência e da linguagem.
Os sentidos participam na construção de \textbf{uma} realidade perceptiva do ser humano, conjuntamente a conceitos, princípios, valores e características individuais \textbf{que} determinam e \textbf{qualificam} as relações das pessoas com o mundo.
Assim, a consciência, então, é formada por processos objetivos e subjetivos de apreensão da realidade, interna e externa, no qual a construção de \textbf{uma} identidade é percebida tanto em termos individuais como no âmbito social.
Normalmente o ser humano se conhece em termos comparativos, ou seja, “ eu sou ” em relação ao outro/outros ou às coisas \textbf{que} circundam.
Barreiro ( 2000, \textbf{pp}. 55-56 ) assevera : > A gênese da consciência humana pode ser explicada pela possibilidade de sua emergência a \textbf{um} nível de descoberta do mundo objetivo como “ forma de realidade ” oposta a \textbf{um} eu individual, mas capaz de integrar -se : consciência-do-mundo x consciência-de-si.
Nessa consciência-do-mundo o \textbf{homem} inclui – na esfera de sua própria dimensão – \textbf{uma} consciência-do-outro, como realidade “ destacada de mim ”, e colocada inicialmente em \textbf{um} mesmo nível. (... ) Para saber -se de si o \textbf{homem} distingue o \textbf{que} está fora ou além deste “ si ” como \textbf{uma} realidade objetiva e intencionalizada por sua consciência.
Essa consciência de existentes-além-de-si contém a própria possibilidade humana de compreender \textbf{uma} \textbf{totalidade} de circunstância em \textbf{que} se colocam : o \textbf{sujeito} consciente ( como o si de sua consciência ) ; o mundo como realidade destacada ( em oposição ) ; o outro ( como outra consciência ) ; a própria circunstância em \textbf{que} se encontram as consciências de \textbf{um} mundo mediatizador.
Barreiro defende a ideia de \textbf{que} existe \textbf{um} processo pelo qual as consciências se reconhecem entre si, conhecem o mundo “ ou adquirem \textbf{uma} visão de mundo ”.
Há, então, \textbf{um} processo de formação do indivíduo, a individualização, \textbf{que} é a tomada da consciência de si em relação aos processos objetivos e subjetivos \textbf{que} ocorrem com cada ser.
Há \textbf{um} processo de aprofundamento quando o indivíduo se \textbf{vê} como \textbf{um} ser social.
Barreiro ( 2000, p.56 ) afirma : > O \textbf{homem} transforma o mundo humanizando -o e transforma o seu próprio modo-de-ser-no-mundo, humanizando -se.
Desde a sua dimensão de Ser histórico o \textbf{homem} tende a se colocar como sujeito de suas relações com o mundo.
Mas, na prática de suas interações, alguns \textbf{homens} podem alienar -se ou ser alienados de \textbf{uma} liberdade geneticamente \textbf{inerente} a eles, mas \textbf{historicamente} inegável, e transformar -se em objetos-para-os-outros ( situação da dominação de consciências ).
Pode -se \textbf{dizer} \textbf{que} se passa para \textbf{um} estado ‘ avançado ’ de consciência quando se coloca crítico em relação à realidade.
Pois mesmo inseridos em \textbf{um} contexto específico se é capaz de olhar com ‘ distanciamento ’ este mesmo espaço/tempo em \textbf{que} se encontra.
Para isso é necessário, em \textbf{um} primeiro momento, saber, conhecer, estar ciente da realidade estabelecida além da \textbf{fronteira} individual.
Para isso o trabalho de conscientização.
Quando alguém se \textbf{vê} no mundo se dá conta de \textbf{que} valores sociais, culturais, econômicos e políticos fazem parte da realidade e de \textbf{que} se está diretamente conectado a eles ou não ( \textbf{depende} da consciência ).
Forma -se, então, a concepção de mundo.
A educação adentra essa realidade como o processo de transmissão de \textbf{um} saber \textbf{que} instrumentaliza o indivíduo na sua relação com esse mundo \textbf{que} está concebendo ou percebendo.
Mesquita ( 2002, p. 79 ) afirma \textbf{que} “ o saber, diferentemente da ciência, \textbf{que} é objetiva, possui a dimensão da subjetividade, da \textbf{cognição}, do pensar sobre o cotidiano, e representa diferentes concepções de mundo ”.
Mais \textbf{adiante} ela enfatiza \textbf{que} “ o saber \textbf{popular} se gesta na vida quotidiana ” : > A vida quotidiana é a vida do \textbf{homem} inteiro, ou seja, o \textbf{homem} participa na vida quotidiana com todos os aspectos de sua individualidade, de sua personalidade.
Nela se colocam ‘ em funcionamento ’ todos os seus sentidos, todas as suas capacidades intelectuais, suas habilidades manipulativas, seus sentimentos, paixões, ideias, ideologias ( HELLER apud MESQUITA, 2002, p.80 ).
Assim, a educação deve ter consciência da subjetividade e habilidades/capacidades contidas em cada indivíduo, \textbf{que} constrói \textbf{uma} diversidade no mundo social ( formado a partir de indivíduos ) na qual ela está a serviço, para poder, assim, auxiliar no desenvolvimento de cada \textbf{um}, em sua inteireza e, consequentemente, da sociedade da qual fazem parte.
O Teatro, assim como as outras expressões da arte, trabalha com instantes significativos da experiência humana, em \textbf{que} a memória é construída a partir dos sentimentos, pensamentos, sensações e ações \textbf{que} produzimos.
Na relação com o objeto artístico, nosso olhar é retribuído, pode tocar o \textbf{íntimo}, trazendo à tona experiências \textbf{cruciais} \textbf{que} tivemos ou \textbf{que} nos motivam novas reflexões.
O encontro com a arte se coloca, desde então, fundamentalmente vinculado com a proposição e a produção de experiências, \textbf{que} nos formam \textbf{enquanto} seres \textbf{humanos}.
Por \textbf{que} pensar o Teatro como instrumento de conscientização?
O Teatro acompanha o ser humano desde os períodos mais remotos da história.
É \textbf{um} fenômeno \textbf{que} faz parte do complexo sistema de vida humana.
E olha o ser humano em sua complexidade. \textbf{Uma} arte \textbf{que} \textbf{exige} muito de quem a pratica e muito mais oferece a quem a aprecia.
Alimenta -se da indelével necessidade de compreensão e ratificação das ideias, sensações, conhecimentos e sentimentos humanos. \textbf{Uma} linguagem \textbf{que} provoca, sensibiliza, emociona.
O Teatro \textbf{enquanto} expressão artística é capaz de gerar transformações no meio social.
Por conta do seu caráter lúdico e a propositura do estado de jogo, torna -se \textbf{um} elo fundamental nos processos de \textbf{ensino-aprendizagem} dos indivíduos, parte importante para \textbf{um} dos objetivos essenciais da educação : o desenvolvimento humano.
Seu essencial instrumento é o próprio ser — a pessoa.
O corpo, a fala, a expressão, a consciência de si mesmo, tudo é \textbf{preciso} ser desenvolvido.
Assim, \textbf{um} trabalho com o Teatro tem início a partir do corpo físico e todas as suas possibilidades, é ele \textbf{que} irá atuar não só no palco, mas em sua vida.
É o seu veículo, \textbf{que} lhe possibilitará andar pelo mundo.
O olhar volta -se para o concreto de si mesmo - a densidade do corpo, sua materialidade, como se movimenta, a função das articulações, o contato com o chão onde os pés pisam e as possibilidades desses braços \textbf{que} estão soltos e podem gesticular.
Toma -se consciência de \textbf{que} se preenche \textbf{um} espaço, interior e exterior.
Pode correr, andar, deitar, sentar e quando nos movimentamos a respiração se \textbf{configura} de forma diferente, pode acelerar, assim como as batidas do coração.
Há \textbf{um} organismo \textbf{que} permite \textbf{que} tudo aconteça e, do ar, através de ondas, pode -se emitir sons ( fala, grita, sussurra, canta, silencia, dialoga ) - descobre -se \textbf{uma} linguagem, \textbf{uma} forma de comunicação.
Toma -se consciência dos sentidos, de emoções e pensamentos.
E os gestos podem expressar o \textbf{que} se pensa e se sente.
O Teatro permite a vivência de experiências sensoriais diversas.
A partir dos jogos teatrais e dramáticos, o jogador pode se posicionar em estado de simulação, envolvendo, nessa realidade paralela do jogo, suas emoções, conceitos, valores.
Japiassu ( 2001, p.26 ) afirma \textbf{que} : > A finalidade do jogo teatral na educação escolar é o crescimento pessoal e o desenvolvimento cultural dos jogadores por meio do domínio, da comunicação e do uso interativo da linguagem teatral, numa perspectiva improvisacional ou lúdica.
O princípio do jogo teatral é o mesmo da improvisação teatral, ou seja, a comunicação \textbf{que} emerge da espontaneidade das interações entre \textbf{sujeitos} engajados na solução cênica de \textbf{um} problema de atuação.
As ações podem ser físicas e/ou psicológicas ( internas ).
A observação de si mesmo, do outro/outra, do espaço externo, do tempo em \textbf{que} se \textbf{vive}, da sociedade em \textbf{que} se está inserido é \textbf{bastante} estimulada no trabalho com o Teatro.
É \textbf{um} instrumento importante para o processo de conscientização e determinação de \textbf{uma} identidade. \textbf{Um} dos elementos fundamentais do Teatro é a ação.
A própria palavra drama, etimologicamente, significa ação.
Interagir, intervir, experimentar, são ações contidas no fazer teatral e permitem \textbf{que} o aluno-jogador se posicione de forma crítica e ativa no processo artístico.
O pensamento ideológico e a formação do senso crítico são condicionados pelos processos criativos, \textbf{vividos}, testados.
Spolin ( 1987, p.03 ) afirma \textbf{que} : “ Se o ambiente permitir, pode -se aprender qualquer coisa, e se o indivíduo permitir, o ambiente lhe ensinará tudo o \textbf{que} ele tem para ensinar. ‘ Talento ’ ou ‘ falta de talento ’ tem muito pouco a \textbf{ver} com isso. ” No Teatro não existe dicotomia entre Arte e conhecimento, como afirma Lima ( 2008, p. 23 ) : > Todas as formas de arte são conhecimento \textbf{sistematizado}, assim como todas as \textbf{categorias} de ciências.
Não há dicotomia entre artes e ciências, elas têm entre si vários pontos \textbf{que} as \textbf{aproximam} e na formação escolar é importante \textbf{que} ambas sejam igualmente valorizadas no currículo.
É necessário superar, também, a concepção de \textbf{que} o conhecimento seja apenas informações.
O conhecimento resulta da ‘ organização ’ de informações em redes de significados.
Teatro é expressão.
Expressar emoções, pensamentos, ideias, ideologias, sua conexão com a vida, a natureza, o divino, \textbf{consigo} mesmo, com o outro, com a sociedade ou o poder estabelecido, algo ele tem para ser \textbf{dito}.
Ele pode mostrar a moralidade de \textbf{uma} situação ou de toda \textbf{uma} época.
Pode rever o passado, compreender o presente ou imaginar o futuro.
A História do Teatro no Mundo nos conta sobre cada civilização \textbf{que} \textbf{viveu} na Terra, seus costumes e mentalidades, assim, vários deixaram seus nomes marcados como grandes dramaturgos, diretores, \textbf{atores}, criadores de \textbf{métodos}, de técnicas, de sistemas.
No Brasil, temos relatos mais claros da existência do teatro desde \textbf{que} os \textbf{jesuítas} chegaram ao país, mas sabe -se \textbf{que} os indígenas \textbf{que} \textbf{viviam} nessas terras utilizavam o teatro como manifestação, inclusive com a utilização de máscaras, danças, cantos e dramatizações.
Nesse sentido, nomes podem ser destacados, entre encenadores, dramaturgos e \textbf{atores}, sem esquecer de outras funções \textbf{que} o teatro abrange, como o iluminador, o figurinista, o cenógrafo, o contrarregra.
Grupos também ficaram \textbf{conhecidos}, principalmente nas últimas décadas, marcando os acontecimentos do \textbf{século} XX e continuam atuantes nos trazendo a realidade e a poesia existente no \textbf{século} XXI.
Falando dos dias \textbf{atuais} em nosso estado, em várias cidades paraibanas, como Guarabira, Areia, Campina Grande, Pedras de Fogo, Alagoa Grande, Bananeiras, Cajazeiras, Conde, Cuité, Monteiro, Souza, Pombal e em sua capital, João Pessoa, acontecem espetáculos em teatros, praças públicas, auditórios e espaços improvisados.
E o teatro é também utilizado por organizações governamentais como a FUNAD ( Fundação Centro Integrado de Apoio ao Portador de Deficiência ) ou não governamentais como a Beira da Linha, ou a Escola Piollin, como \textbf{um} instrumento de reconhecimento da identidade e de processo de conscientização — do indivíduo inserido em \textbf{um} contexto regional e universal.
Em todas as experiências, os resultados relatados são rápidos e profundos.
O protagonismo encontrado no teatro é e pode ser melhor utilizado na educação como caminho para \textbf{que} todos reconheçam \textbf{que} são heróis de sua própria história e não as suas vítimas.
E \textbf{que} podemos e devemos ser cidadãos do mundo, seres humanos \textbf{que} sentem, pensam, falam e agem.
Estamos \textbf{aqui} para realizar o grande e maravilhoso propósito da vida, a grande obra de arte \textbf{que} o teatro \textbf{revela}.
A linguagem do Teatro, especificamente no âmbito do Ensino Médio, está assegurada legalmente, tendo em vista a fundamentação da concepção de Educação expressa nos diversos dispositivos \textbf{que} a regulamentam, assegurando aos(às) estudantes, na esfera da educação formal, o alcance do objetivo proposto pela Base Nacional Comum Curricular para o Ensino Médio ( 2018 ).
Através da definição de competências e habilidades, a BNCCEM se propõe como objetivo “ consolidar, aprofundar e ampliar a formação integral dos estudantes, atendendo às finalidades dessa etapa e contribuindo para \textbf{que} cada \textbf{um} deles possa construir e realizar seus projetos de vida, em consonância com os princípios da justiça, da ética e da cidadania ”. ( BNCCEM, 2018, p. 470 ) Para garantir o cumprimento dos objetivos propostos pela BNCCEM, através das Habilidades e Competências para cada linguagem e, portanto, assegurando o direito de aprendizagem dos/as estudantes, faz -se \textbf{imprescindível} a atuação de profissionais capacitados e habilitados em sua área específica, a saber, \textbf{um} professor licenciado em teatro ou seu equivalente. \textbf{Uma} vez garantida a atuação profissional adequada e suficiente, através da implementação das mudanças decorrentes da Lei 13278/2016, pode -se vislumbrar o desenvolvimento de \textbf{um} processo de aprendizagem do Teatro e consonante com os princípios de \textbf{uma} educação \textbf{que} visa “ a formação e o desenvolvimento humano global, em suas dimensões intelectual, física, afetiva, social, ética, moral e simbólica ”. ( BNCCEM, 2018, p. 16 ) O ensino da linguagem teatral no Ensino Médio deve estar articulado às aprendizagens estabelecidas para o Ensino Fundamental, de modo \textbf{que}, neste percurso, sejam aprofundados, ampliados e consolidados os conhecimentos próprios do Teatro.
Neste sentido, é indispensável abalizar, de \textbf{um} \textbf{lado}, os atributos específicos da área teatral e, de outro, indicar as dimensões de aprofundamento a serem desenvolvidas no Ensino Médio.
De tal modo, \textbf{pontua} -se \textbf{que} o ensino do Teatro, assim como as demais linguagens, tem \textbf{fundamentos} próprios e \textbf{abordagens} pedagógicas distintas e \textbf{que}, no âmbito do ensino formal, não pode ser confundido com \textbf{mera} atividade de lazer ou recreação.
O ensino de Teatro tem por pressupostos o desenvolvimento das capacidades do indivíduo de se (re)conhecer \textbf{enquanto} sujeito, de se relacionar com o outro, de reconhecer estéticas e discursos \textbf{historicamente} construídos e transformados na relação espaço-temporal.
Cada \textbf{um} destes pressupostos \textbf{implica} em \textbf{desdobramentos} relativos tanto à práxis teatral quanto ao desenvolvimento pleno do \textbf{sujeito}.
Desse modo, discorrer sobre a capacidade do \textbf{sujeito} de se reconhecer, na perspectiva teatral, \textbf{implica} em \textbf{um} contínuo trabalho sobre o próprio corpo ( corpo-mente ) e compreensão de si na esfera da criação cênica ; de outro \textbf{lado}, essa dimensão igualmente precisa ser ponderada na relação do \textbf{sujeito} com a esfera social, econômica e cultural em \textbf{que} está inserido, \textbf{uma} vez \textbf{que} estes parâmetros balizam a formação do \textbf{sujeito}.
Como consequência, pode -se refletir acerca das implicações na esfera de relação com o outro, ou seja, à medida \textbf{que} \textbf{um} \textbf{sujeito} compreende seu processo e formação, também pode assumir \textbf{um} olhar sensível para com o outro.
Igualmente, pode -se sinalizar as possíveis interconexões entre a dimensão estética e a produção teatral \textbf{historicamente} construída, bem como, a possibilidade de reflexão sobre os elementos \textbf{que} compõem o fazer teatral ( dramaturgia, cenografia, atuação, iluminação, figurinos, sonoplastia, etc. ) em diferentes períodos e nas diferentes sociedades.
Estas são apenas algumas das possibilidades para o ensino de Teatro \textbf{que}, todavia, não pode prescindir de seu \textbf{fundamento} : a relação entre “ atores/atrizes ”.
Atores/atrizes, neste sentido, se refere aos protagonistas no processo de \textbf{ensino-aprendizagem} : os estudantes e professores.
O teatro, \textbf{enquanto} arte coletiva e efêmera ( ainda \textbf{que} passível de registro ) se funda na relação entre os \textbf{sujeitos}, \textbf{ora} como atuantes, \textbf{ora} como espectadores/as.
Na \textbf{abordagem} do jogo teatral, por exemplo, esta relação é evidenciada por Ingrid Koudela ( 2011, p. 241, \textbf{grifos} do \textbf{autor} ) \textbf{que}, pautada na \textbf{abordagem} triangular de Ana Mae Barbosa, compartilha o \textbf{fundamento} do trabalho com Teatro no âmbito educacional : > O jogo teatral é \textbf{um} jogo de construção em \textbf{que} a consciência do ‘ como se ’ é gradativamente trabalhada em \textbf{direção} à articulação da linguagem artística do teatro.
No processo de construção dessa linguagem, a criança e o jovem estabelecem com seus pares \textbf{uma} relação de trabalho, combinando a imaginação dramática com a prática e a consciência na observação das regras do jogo teatral.
Neste sentido, cabe ressaltar \textbf{que}, em se tratando do Ensino Médio e consoante com a proposta de construção e realização de projetos de vida pelos estudantes, é pertinente aprofundar os conhecimentos na área teatral, ainda \textbf{que} o objetivo não seja o de transformar os/as estudantes em “ artistas ”, mas sim, de oferecer -lhes \textbf{uma} base de conhecimento \textbf{sólida} em relação às esferas de atuação no âmbito do Teatro, no cotejamento com a formação cidadã do \textbf{sujeito} e o entendimento da dimensão humana da vida.
Em relação à faixa etária em \textbf{que} normalmente os estudantes adentram ao Ensino Médio, a adolescência corresponde à fase da vida em \textbf{que} o indivíduo muitas vezes “ se apresenta como criança em seus jogos, brincadeiras e como adulto com seu corpo, seus novos sentimentos e suas expectativas de futuro. ” ( COSTA ; MEDEIROS ; \textbf{SANTOS}, 2011, p.251 ) Há \textbf{um} processo biopsicossocial do indivíduo \textbf{que} ocorre, trazendo necessidade de pertencimento à \textbf{um} grupo e sentimentos \textbf{que} precisam ser reprimidos ou escondidos, como a insegurança ou o “ não saber sobre alguma coisa ”.
Deve considerar -se \textbf{que} as \textbf{influências} de múltiplas variáveis sobre o comportamento dos adolescentes, incluindo as autopercepções, alteram o senso de autoeficácia, senso este \textbf{que} nos permite conhecer as reais potencialidades e os limites, buscando de forma mais concisa o desenvolvimento de novas habilidades \textbf{que} levarão o indivíduo as conquistas as quais ele almeja. ( COSTA, MEDEIROS E \textbf{SANTOS}, 2011, p.252 ).
Nesse processo de mudanças emocionais, psíquicas e biológicas, \textbf{que} se refletem no comportamento social, a escola possui papel importante como espaço de construção e desenvolvimento de habilidades \textbf{que} auxiliam na autoestima e auto \textbf{percepção}.
É \textbf{preciso} cuidar para \textbf{que} a questão da “ eficácia ” colocada por Costa, Medeiros e Santos, além de outros \textbf{pesquisadores} da atualidade, seja observada com base \textbf{que} contemple fatores objetivos e subjetivos próprios da \textbf{humanidade}.
Pessoas respondem de forma diferente aos \textbf{estímulos} internos e externos.
Assim, quando pensamos em processo pedagógico, precisamos ter em \textbf{mente} \textbf{que} estudantes apreendem e percebem o mundo de forma diversa e \textbf{que} a inteligência não pode ser compreendida apenas por \textbf{um} aspecto intelectual, mas a partir da complexidade \textbf{que} nos é \textbf{inerente}.
Habilidades nos mostram dimensões cognitivas, emocionais, psicológicas e sociais \textbf{que} são colocadas em ação conforme o indivíduo \textbf{vive} as experiências na vida. \textbf{Um} dos grandes objetivos da Educação é \textbf{que} o estudante se responsabilize pelo próprio aprendizado e, para isso, a escola \textbf{necessita} de metodologias e espaços para \textbf{que} o indivíduo desenvolva de forma integrada e contínua os aspectos \textbf{que} lhe compõe.
O Teatro faz parte do processo pedagógico no Ensino Médio como \textbf{um} espaço possível de união da \textbf{teoria} e da prática, através de exercícios, jogos, histórias, brincadeiras, diálogos, encenações, improvisações, trabalho físico, trabalho vocal, relações simbólicas, leitura e interpretação de textos, desenvolvimento dos sentidos, reconhecimento do espaço, do tempo, das emoções, dos processos mentais, a escrita, a sociabilização, etc.
Além da utilização de todas as outras expressões artísticas, como as Artes Visuais, a Música e a Dança.
O Teatro nas escolas está comprometido com o desenvolvimento humano, com indicativos organizados e de acordo com cada ano do Ensino Médio.

% matched lemmas: abordagem, adiante, aproximar, aqui, ator, atual, autor, bastante, categoria, cognição, configurar, conhecido, conseguir, crucial, depender, desdobramento, direção, dizer, enquanto, ensino-aprendizagem, estímulo, exigir, fronteira, fundamento, grifo, historicamente, homem, humanidade, humanos, implicar, imprescindível, inerente, influência, jesuíta, lado, mente, mero, moderno, método, necessitar, ora, percepção, pesquisador, pontuar, popular, pp, preciso, qualificar, que, revelar, santo, sistematizar, sujeito, século, sólido, teoria, totalidade, um, ver, viver, íntimo
\end{textsample}
